\section{Theoretical Analysis}\label{sec:theory}

This section provides four formal results that anchor the empirical findings of Section \ref{sec:experiments}: a counting argument that quantifies the parameter-matching condition (Proposition \ref{prop:param}), the resulting communication-cost ratio (Proposition \ref{prop:comm}), a B-spline approximation bound specialised to NetFlow features (Lemma \ref{lem:spline}), and a convergence theorem for FedKAN under non-IID partitioning (Theorem \ref{thm:conv}). Together these results predict that the heterogeneity-amplified variance term $\Gamma_\alpha$ couples to a constant spline-bias floor $\mathcal{O}(G^{-(k+1)})$ rather than to the architectural parameters of the network, a prediction we test in Section~\ref{sec:experiments} and do not confirm: the seed-variance ratio it would imply holds on one dataset and reverses on the other three.

\subsection{Parameter Complexity}

\begin{proposition}[Parameter Complexity of KAN vs MLP]\label{prop:param}
A KAN layer (Definition \ref{def:kan}) with input dim $n_{\text{in}}$ and output dim $n_{\text{out}}$, grid size $G$, and order $k$ has
\begin{equation}\label{eq:kan-count}
|\Theta_{\text{KAN}}| \;=\; n_{\text{in}}\, n_{\text{out}} \,(G + k + 2),
\end{equation}
while an MLP layer of the same dimensions with bias has
\begin{equation}\label{eq:mlp-count}
|\Theta_{\text{MLP}}| \;=\; n_{\text{in}}\, n_{\text{out}} + n_{\text{out}}.
\end{equation}
\end{proposition}

\begin{proof}
By Definition \ref{def:kan}, each edge $\phi_{q,p}$ contributes one base scalar $w^{(b)}_{q,p}$, one spline scalar $w^{(s)}_{q,p}$, and $G+k$ B-spline coefficients $c_{q,p,i}$, totalling $G+k+2$ parameters. There are $n_{\text{in}}\cdot n_{\text{out}}$ edges, giving \eqref{eq:kan-count}. The MLP layer has weight matrix $W \in \mathbb{R}^{n_{\text{out}} \times n_{\text{in}}}$ and bias $\mathbf{b} \in \mathbb{R}^{n_{\text{out}}}$, summing to \eqref{eq:mlp-count}.
\end{proof}

\begin{corollary}[Parameter-Matched Hidden Width]\label{cor:match}
For a single-hidden-layer architecture $n_{\text{in}} \to n_h \to n_{\text{out}}$, parameter parity between KAN (hidden width $n_h^{\text{KAN}}$, hyperparameters $G, k$) and MLP (hidden width $n_h^{\text{MLP}}$) requires, to leading order in the dominant linear term,
\begin{equation}
n_h^{\text{MLP}} \;\approx\; (G + k + 2)\, n_h^{\text{KAN}}.
\end{equation}
For our experimental configuration ($G=5$, $k=3$) the matching ratio is $10$, motivating the comparison between KAN-8 ($\sim 3.3$k parameters) and MLP-PM-80 ($\sim 3.4$k parameters) used in Section \ref{sec:experiments}.
\end{corollary}

\subsection{Communication Complexity}

\begin{proposition}[Communication Complexity Ratio]\label{prop:comm}
Under FedAvg with $T$ rounds, fraction $C \in (0, 1]$ of $K$ clients per round, and float32 transmission ($4$ bytes per scalar), the total uplink+downlink bytes for an architecture with $|\Theta|$ parameters is
\begin{equation}
\mathcal{C}_{\text{total}}(|\Theta|) \;=\; 8\, T\, C\, K\, |\Theta|.
\end{equation}
For a single-hidden-layer architecture, the KAN/MLP ratio simplifies to
\begin{equation}\label{eq:comm-ratio}
\frac{\mathcal{C}_{\text{KAN}}}{\mathcal{C}_{\text{MLP}}} \;=\; \frac{(G+k+2)\, n_h^{\text{KAN}}\,(n_{\text{in}}+n_{\text{out}})}{(n_{\text{in}}+1)\, n_h^{\text{MLP}}\, n_{\text{out}} + n_h^{\text{MLP}}}.
\end{equation}
\end{proposition}

\noindent The ratio in \eqref{eq:comm-ratio} clarifies a misconception: KAN architectures are
\emph{not} intrinsically communication-efficient. They are efficient only when the chosen
$n_h^{\text{KAN}}$ is much smaller than the MLP's at matched accuracy. Under parameter parity
the KAN's per-round uplink exceeds the MLP's.

\textbf{A correction to this proposition, found while responding to review.} As stated,
\eqref{eq:comm-ratio} counts learnable parameters, and for the MLP that is exactly what is
transmitted: the measured uplink matches the formula to within $0.0\%$ for both MLP widths.
For the KAN it does not. FedKAN-8 transmits $153{,}760$~B per round across ten clients where
the proposition predicts $131{,}200$~B, an excess of $17.2\%$; FedKAN-16 shows $10.1\%$. The
excess is not noise and not an implementation defect: a KAN layer's \texttt{state\_dict}
carries its B-spline knot vector as a buffer, $468$ and $96$ scalars for the two layers of
FedKAN-8, which is $564$ non-learnable values against $3{,}280$ learnable ones, or exactly the
$17.2\%$ observed. Any implementation that synchronises \texttt{state\_dict} rather than
\texttt{parameters()} pays this, and ours does.

We therefore restate the communication accounting in terms of transmitted scalars rather than
learnable parameters. The qualitative conclusion of the first version, that under parameter
parity the KAN transmits more, survives; the mechanism we gave for it did not include the knot
buffers and so the figure we quoted was right for an incomplete reason.

\subsection{Spline Approximation on NetFlow Features}

\begin{lemma}[B-spline Approximation Bound]\label{lem:spline}
Let feature $j$ be supported on $[a_j, b_j]$ with density $\rho_j$ that is $L_j$-Lipschitz, and let $\hat\phi_j$ denote the order-$k$ B-spline approximation of any target $\phi_j \in C^{k+1}([a_j, b_j])$ on a uniform grid of $G$ points. Then
\begin{equation}\label{eq:spline-bound}
\|\phi_j - \hat\phi_j\|_\infty \;\leq\; \frac{L_j (b_j - a_j)^{k+1}}{(k+1)!\, G^{k+1}}.
\end{equation}
\end{lemma}

\begin{proof}
Standard B-spline approximation result; see \cite[Theorem 6.27]{schumaker2007}.
\end{proof}

\textbf{Implication for NetFlow.} Standardised NetFlow-v2 features \cite{sarhan2022standard} after Min-Max scaling lie in $[0, 1]$, and many flow-statistic features (e.g. \texttt{FLOW\_DURATION}, \texttt{TCP\_FLAGS}, \texttt{IN/OUT\_BYTES}) exhibit smooth empirical densities once log-transformed. The bound \eqref{eq:spline-bound} therefore predicts a small approximation residual at modest grid sizes, motivating our choice of $G=5, k=3$.

\subsection{Convergence under Non-IID}

\begin{theorem}[FedKAN Convergence under Heterogeneity]\label{thm:conv}
Suppose each client loss $\mathcal{L}_k$ is $L$-smooth and $\mu$-strongly convex in $\Theta$, the per-client gradient variance is bounded above by a heterogeneity factor $\Gamma_\alpha = \mathbb{E}\,\|\nabla \mathcal{L}_k(\Theta) - \nabla F(\Theta)\|^2$ that is monotonically decreasing in the Dirichlet concentration $\alpha$, and a diminishing learning-rate schedule satisfying the conditions of \cite{li2020convergence}. Let $\Theta_T$ be the FedKAN-IDS global iterate after $T$ communication rounds. Then
\begin{equation}\label{eq:thm}
\mathbb{E}\big[F(\Theta_T) - F(\Theta^\star)\big] \;\leq\; \frac{2L}{\mu^2 T}\big(L + \Gamma_\alpha\big) \;+\; \mathcal{O}\!\left(G^{-(k+1)}\right),
\end{equation}
where the additive $\mathcal{O}(G^{-(k+1)})$ residual is the spline-approximation bias inherited from Lemma \ref{lem:spline} aggregated across all KAN edges, and is independent of the heterogeneity factor $\Gamma_\alpha$.
\end{theorem}

\begin{proof}[Sketch]
The dominant non-IID FedAvg term $\frac{2L}{\mu^2 T}(L+\Gamma_\alpha)$ is obtained by adapting Theorem 2 of \cite{li2020convergence} to the parameter set $\Theta_{\text{KAN}}$, which remains a finite-dimensional Euclidean space and inherits the smoothness/strong-convexity assumptions because the SiLU base activation and linear B-spline expansion in \eqref{eq:edge} are differentiable. The spline-bias term arises because each edge function $\phi_{q,p}$ is approximated only up to the bound in Lemma \ref{lem:spline}; aggregating the per-edge $L_\infty$ errors over $n_{\text{in}}\cdot n_{\text{out}}$ edges yields an overall residual of order $G^{-(k+1)}$ in the loss, which does not couple to the gradient-heterogeneity term $\Gamma_\alpha$.
\end{proof}

\textbf{An empirical check on the assumptions.} Reviewer feedback asked for evidence on the
local convexity that Theorem~\ref{thm:conv} assumes, and the evidence is unfavourable. Running
Lanczos with full re-orthogonalisation on the exact Hessian-vector product at the converged
federated solution ~\cite{ghorbani2019hessian,yao2020pyhessian} (NF-BoT-IoT-v2, binary, $50$ rounds, $8192$ held-out samples, $40$ Ritz
directions) gives $\lambda_{\max}=88.5$ and $\lambda_{\min}=-0.134$ for FedKAN-8, with $45\%$
of the Ritz values negative, and $\lambda_{\max}=0.52$, $\lambda_{\min}=-1.8\times 10^{-3}$
for the parameter-matched MLP, with $27.5\%$ negative. Negative curvature at the solution means
strong convexity fails locally as well as globally, so the theorem describes a regime the
training does not reach. We keep it for the structure it exposes, that the spline bias enters
additively and does not scale with $\Gamma_\alpha$, and not as a bound on our runs. The two
condition numbers are the same order ($662$ against $285$), so the spectrum does not explain
any performance difference between the architectures either.

% GO BO 19/08/2026. He qua nay KHONG suy ra duoc tu Dinh ly 1: so hang
% 2L/mu^2 T (L + Gamma_alpha) ap giong het cho ca KAN lan MLP, con sai so spline
% O(G^{-(k+1)}) la sai so CONG THEM rieng cho KAN chu khong phai la chan. Dao ham
% theo Gamma_alpha bang nhau o hai kien truc, nen ket luan "KAN it nhay hon voi
% khong dong nhat" la mot KHANG DINH THEM. Phan thuc nghiem sau khi do lr rieng
% cho tung kien truc cung khong con ung ho ket luan do (ty le phuong sai 2,62x -> 1,13x).
% Nguyen van he qua cu duoc giu o FINDING-R2-7-learning-rate.md.

